% ============================================================
% Section 5: Limitations  +  Section 6: Conclusion
% Overleaf 直贴版（Springer LNCS llncs 模板）
% 用户五项决定：1B 离线精排缺陷不写入；2A 红旗 70% 写入；
%               3B 外部基线缺失不明说；4B 未来工作两条；5 标题 Limitations
% 数字均对照 v1 归档 pipeline_eval_v1_30_final.json 与 data_pipeline_technical_report.md
% ============================================================

\section{Limitations}\label{sec:limits}

The current system and evaluation have the following limitations.

\textit{Data.} All 1{,}339 records of CMCR-Colloquial are currently in draft
status, and the double-blind expert review has not yet been performed. The
colloquial consultations and triage outputs were generated by a single LLM
(DeepSeek-V4.1-Flash), and the proposition decomposition was produced by the
same model; the automated quality-control scores (five dimensions on each side
of the pipeline, out of 10) therefore provide only a lower-bound quality
guarantee rather than a discriminative quality measure. The dataset currently
covers Chinese only, and its transferability to other languages and healthcare
settings is unverified.

\textit{Evaluation.} The end-to-end evaluation uses 30 samples, and the judge
shares the model pool with the generator, so self-preference bias cannot be
excluded and the mean score of 71.4 should be read as a relative reference.
The immediate-emergency-care recommendation is issued for 70\% of red-flag
queries; the remaining three of ten were not escalated. This recall gap
motivated the recall-oriented red-flag rescue mechanism of
Sect.~\ref{sec:architecture}, which prefers escalation over omission; the
mechanism was added after the evaluation, and its effect has not yet been
systematically evaluated. The LoRA-based dense retriever is a negative result
(mean Recall@50 of 0.197 over five folds), and no corrective hyperparameter
search has been attempted.

\textit{System.} The end-to-end latency and API cost of Jev verification and
generation have not yet been systematically profiled.

\section{Conclusion}\label{sec:conc}

We presented FirstCare, a low-hallucination Chinese triage demonstration
system driven by a real-world case data resource and dual-question fact
verification: 1{,}543 peer-reviewed clinical case reports are structured and
decomposed into 43{,}271 atomic propositions as a precomputed verification
base, and 1{,}339 colloquial question--answering pairs support the end-to-end
demonstration. The dual-question protocol partitions patient statements into
three states; \textit{absent} propositions are removed by a hard gate,
red-flag propositions trigger mandatory escalation, and a two-column split
together with a two-layer post-hoc fact guard constrains generation. The
system runs offline in a single-file web application and is bit-identical to
the Python backend. Future work includes completing the double-blind expert
review and releasing the dataset, and conducting larger-scale human expert
evaluation with external baselines, systematically evaluating the red-flag
rescue mechanism, and profiling the cost of verification and generation at
scale.

% ============================================================
% 提交前待办（全文组装清单）：
% 1. 统一系统名 FirstCare（标题/摘要/正文）；
% 2. Keywords 填写；\titlerunning{} 设置（消除第 3 页页眉报错）；
% 3. demo URL 脚注替换为真实托管地址（Sect.4 占位）；
% 4. 清除模板残留（"150–250 words."、"1.1 A Subsection Sample"、贡献3"Legal-Aid"）；
% 5. Intro 压缩 ~0.7 页；Fig.1/Fig.2 PNG 拖入 Overleaf 工程；
% 6. 评测表若上稿须从 v1 归档重新导出（export_latex_tables.py 当前为 v2）；
% 7. 参考文献插入（CMCR 条目见 Sect.2 文件注释）。
% ============================================================
